\documentclass[11pt]{article}
\usepackage[utf8]{inputenc}
\usepackage[T1]{fontenc}
\usepackage{graphicx}
\usepackage{amsmath}
\usepackage{amssymb}
\author{Hesam Pakdaman}
\date{\today}
\title{}


\newcommand{\dist}[3]{d_#1(#2, #3)}
\newcommand{\distfive}[2]{\frac{|#1|}{1 + |#2|}}
\newcommand{\distfivenoabs}[2]{\frac{#1}{1 + #2}}

\begin{document}

\begin{itemize}
\item[\noindent\textbf{11.} $\dist{1}{x}{y}$] This is not a metric
  since 2.15 (c) is not satisfied. We give an example
  $\dist{1}{10}{0} = 100$, $\dist{1}{10}{4} = 36$, and
  $\dist{1}{4}{0} = 16$ so that
  \begin{align*}
    \dist{1}{10}{0} > \dist{1}{10}{4} + \dist{1}{4}{0}.
  \end{align*}

\item[$\dist{2}{x}{y}$] Both 2.15 (a) and (b) are clearly true. We
  show that (c) is also satisfied. Assume not, then there exists
  points $x$, $y$ and $r$ such that
  \begin{align*}
    \dist{2}{x}{y} > \dist{2}{x}{r} + \dist{2}{r}{y},
  \end{align*}
  which in this particular case is
  \begin{align*}
    \sqrt{|x-y|} > \sqrt{|(x-r)|} + \sqrt{|r-y|}.
  \end{align*}
  If $p > q > 0$ then $p^2 > q^2$ for any $p,q \in R$ so that
    \begin{align*}
      |x-y| &> (\sqrt{|(x-r)|} + \sqrt{|r-y|})^2 \\
            &= |x-r| + 2\sqrt{|x-r| |r-y|} + |r-y| \\
            &\geq |x-r| + |r-y|,
  \end{align*}
  where the last inequality comes from the fact that
  $\sqrt{|x-r| |r-y|} \geq~0$. The above shows that $d(p,q) = |p - q|$
  cannot be a metric. But that is a contradiction since Theorem 1.37
  shows that $|p-q|$ satisfies 2.15 (c).

\item[$\dist{3}{x}{y}$] Conditions 2.15 (a), (b) are true and we show
  that (c) is satisfied. For any $r \in R$ we have that
  \begin{align*}
    \dist{3}{x}{y} &= |x^2 - y^2| = |(x^2-r^2) + (r^2-y^2)| \\
                   &\leq |x^2-r^2| + |r^2-y^2| \\
                   &= \dist{3}{x}{r} + \dist{3}{r}{y},
  \end{align*}
  where we have used Theorem 1.37 to get the inequality. This shows
  that $\dist{3}{x}{y}$ is a metric.

\item[$\dist{4}{x}{y}$] We show that this is not a metric since
  condition (a) of Definition 2.15 is not satisfied. If $x \neq 0$
  then
  \begin{align*}
    \dist{4}{x}{x} = |x - 2x| = |x| > 0.
  \end{align*}

\item[$\dist{5}{x}{y}$] Conditions in (a), (b) are satisfied so we
  focus on (c). Throughout this exercise WLOG assume $x \leq y$.

  Suppose $x \leq r \leq y$, we have that
  \begin{align*}
    \dist{5}{x}{y} &= \distfive{x-y}{x-y} = \distfive{x-r+r-y}{x-y} \\
                   &\leq \distfive{x-r}{x-y} + \distfive{r-y}{x-y} \\
                   &\leq \distfive{x-r}{x-r} + \distfive{r-y}{r-y},
  \end{align*}
  where we used Theorem 1.37 in the first inequality. Since
  $x \leq r \leq y$ we know that $|x-r| \leq |x-y|$ and
  $|r-y| \leq |x-y|$, both of which we used to get the second
  inequality.

  Now assume $x \leq y < r$. First we show that for any
  $\varepsilon > 0$
  $$\distfive{w}{w} < \distfive{w+\varepsilon}{w+\varepsilon},$$
  holds. The statement is clearly true for $w=0$. If $|w| > 0$ then we
  can write
  \begin{align} \label{eq:bigger}
    \distfive{w}{w} = \distfivenoabs{1}{\frac{1}{|w|}}
                    < \distfivenoabs{1}{\frac{1}{|w+\varepsilon|}} 
                    = \distfive{w+\varepsilon}{w+\varepsilon},
  \end{align}
  where we get the first inequality since it is true for denominator
  that $ 1+ \frac{1}{|w+\varepsilon|} < 1 + \frac{1}{|w|}$. Since
  $x \leq y < r$ we have that $|x-y| < |x - r|$ which means we can use
  \eqref{eq:bigger} to get the first inequality below
  \begin{align*}
    \distfive{x-y}{x-y} < \distfive{x-r}{x-r} \leq \distfive{x-r}{x-r} + \distfive{r-y}{r-y}.
  \end{align*}
  The last inequality is due to the last term being
  non-negative. Similar argument can be made for $r < x \leq y$
  because $|x-y| < |r-y|$ which allows us to use \eqref{eq:bigger}
  again. This shows that $\dist{5}{x}{y}$ is a metric.


\end{itemize}



\end{document}